% ── Conclusiones ─────────────────────────────────────────────────────────────
\section{Conclusiones}

\subsection{Resultados}
% TODO(agente): Marcar con [Logrado] / [Parcial] / [No logrado] cada objetivo
% específico listado en la Introducción, según lo que esté realmente
% implementado y verificable en el código al momento de generar el informe.
% No declarar un objetivo como logrado si no hay evidencia en el repositorio.
\begin{itemize}
    \item[] Diseño del modelo de base de datos relacional: \_\_\_\_\_\_\_\_\_\_.
    \item[] Implementación de la API REST: \_\_\_\_\_\_\_\_\_\_.
    \item[] Implementación de la interfaz web (asistente de trámite y panel administrativo): \_\_\_\_\_\_\_\_\_\_.
    \item[] Mecanismos de autenticación y autorización: \_\_\_\_\_\_\_\_\_\_.
\end{itemize}

\subsection{Conclusiones}
% TODO(agente): Redactar entre 3 y 5 conclusiones basadas únicamente en lo
% efectivamente construido y verificado, evitando afirmaciones generales no
% respaldadas por el sistema implementado.
\begin{enumerate}
    \item El uso de una arquitectura desacoplada entre el frontend en React
    y el backend en Express permitió desarrollar y probar ambas capas de
    forma independiente, comunicándose únicamente a través de contratos de
    API bien definidos.
    \item El modelo de datos relacional en PostgreSQL, con restricciones
    \texttt{CHECK} sobre los campos de estado, permitió delegar parte de la
    validación del flujo del trámite a la propia base de datos, reduciendo
    la posibilidad de estados inconsistentes.
    \item[] % TODO(agente): agregar conclusiones adicionales específicas del
    % resultado de las pruebas realizadas.
\end{enumerate}

\subsection{Recomendaciones}
\begin{enumerate}
    \item Para una versión posterior al alcance de demo, se recomienda
    incorporar pruebas automatizadas (unitarias y de integración) sobre los
    servicios de la API, particularmente sobre la lógica de cambio de
    estado de una solicitud.
    \item Se recomienda evaluar la integración con un proveedor de pagos
    real y con el sistema de autenticación institucional de la UNSAAC para
    una eventual puesta en producción.
    \item[] % TODO(agente): agregar recomendaciones adicionales pertinentes
    % al estado real del proyecto.
\end{enumerate}
